\section{Research question}
In financial disturbances, the trajectories implied by a recently estimated
smooth trend may diverge markedly from later observations. It remains unclear
whether such divergence indicates informative departures from prevailing
conditions or merely ordinary shifts in forecast uncertainty.
Our question is: across countries, index-construction methods, and forecast
horizons, when do smoothed financial-stress-index forecasts cease to be
reliable, and do errors reveal information beyond conventional models?

\subsection{Scope and contribution}
The intended contribution is a transparent cross-country empirical assessment
rather than a new penalized smoothing formula, numerical optimizer or
dynamic-factor estimator. We plan to study Mexico, the United States,
the United Kingdom, Japan and a defined euro-area aggregate, subject to
audited series availability. Earlier U.S.-only retrospective results motivate,
but do not establish, the proposition.
